\documentclass[10pt]{article}
\usepackage[margin=20mm]{geometry}
\usepackage{amsmath,amssymb,fontspec,unicode-math,xurl,hyperref}
\IfFontExistsTF{DejaVu Serif}{\setmainfont{DejaVu Serif}}{\setmainfont{Latin Modern Roman}}
\setmathfont{latinmodern-math.otf}
\setlength{\parindent}{0pt}\setlength{\parskip}{6pt}\sloppy
\begin{document}
Abstract: The published probabilistic-attention objective differs from the claimed Gaussian negative log-likelihood and reverses its scale incentive in an admissible scalar case. This is a Major concern about the specification of a core module. Benchmark arithmetic and the principal ablation differences are consistent; implementation and empirical effects remain unresolved because source code was unavailable.\par
Author: AI-assisted technical review\par
Date:\par
Title: Technical Verification: S2-Track: A Simple yet Strong Approach for End-to-End 3D Multi-Object Tracking\par
\# Identity and scope\par
S2-Track: A Simple yet Strong Approach for End-to-End 3D Multi-Object Tracking. Tao Tang; Lijun Zhou; Pengkun Hao; Zihang He; Kalok Ho; Shuo Gu; Zhihui Hao; Haiyang Sun; Kun Zhan; Peng Jia; Xianpeng Lang; Xiaodan Liang. PMLR 267 (2025), paper I25-43683. The reviewed source is the published PMLR PDF frozen for B004. Page numbers below are one-based PDF pages (also the article page labels), not volume pagination. This is an initial technical review; scientific adjudication remains pending.\par
The main text, pp. 1–9, and embedded Appendices A–D, pp. 12–17, were read. The focused technical check concerns the probabilistic decoder’s objective, its integration into the full loss, and arithmetic in the comparison and ablation tables. No public author implementation was available in the frozen packet. The supplementary video was unavailable. Thus the implemented historical loss and benchmark scores cannot be established independently here.\par
\# Definitions and the affected objective\par
Section 3.3 models each attention value by a Gaussian and gives the reparameterization \(\alpha=\mu+\sigma\epsilon\), \(\epsilon\sim N(0,1)\), where \(\sigma>0\) is a standard deviation. Write \(m=q_i k_j/\sqrt d\) and \(r=\alpha-m\). The description calls Eq. (6), p. 5, a negative log-likelihood. The actual printed summand, confirmed in the original page image, is \[\ell_{\rm print}(r,\sigma)=\log\left(\frac1{\sqrt{2\pi\sigma^2}}\right)\exp\left(-\frac{r^2}{2\sigma^2}\right).\] There is no leading minus sign, and the logarithm closes before the exponential. Section 3.5 explicitly incorporates this loss into Eq. (10) with a positive coefficient, default one; Appendix A.2.3 repeats that choice.\par
\# Substantive finding\par
Problem 1 (The specified UPD loss reverses the Gaussian scale incentive).  The printed loss is not the claimed Gaussian negative log-likelihood. The difference changes the optimization direction even in a single admissible scalar attention entry. It affects the mathematical specification of one of the paper’s three core modules.\par
\#\#\#\# Complete calculation.\par
Equation reference: eq:nll. For an observation \(\alpha\) from the stated Gaussian centered at \(m\), the negative logarithm of its density is \[
 \ell_{\rm NLL}(r,\sigma)=\log(\sqrt{2\pi}\sigma)+\frac{r^2}{2\sigma^2}.\] Take \(m=\mu=0\), allowed by zero query/key dot product and an MLP output of zero mean, and any positive \(\sigma\). At the fixed residual \(r=0\), the printed expression is \(-\log(\sqrt{2\pi}\sigma)\), with derivative \(-1/\sigma\); the genuine NLL has derivative \(+1/\sigma\). This is not merely an additive constant or harmless normalization.\par
The distinction also survives the reparameterization in the paper, so it is not dependent on a probability-zero draw. With \(\alpha=\sigma\epsilon\), direct Gaussian integration gives \[\begin{aligned}
 \mathbb E\ell_{\rm print}(\sigma\epsilon,\sigma)
 &=-\frac{\log(\sqrt{2\pi}\sigma)}{\sqrt2},\\
 \mathbb E\ell_{\rm NLL}(\sigma\epsilon,\sigma)
 &=\log(\sqrt{2\pi}\sigma)+\frac12.\end{aligned}\] Indeed, \(\mathbb E e^{-\epsilon^2/2}=1/\sqrt2\) and \(\mathbb E\epsilon^2=1\). Thus the printed expected loss decreases without bound as \(\sigma\to\infty\), whereas the stated NLL increases. The reviewer-written script independently integrates the former for \(\sigma=0.1,1,10,100\) and agrees with the closed form within \(10^{-10}\). This proves an error in the specified component objective; it does not prove that the full tracking loss is unbounded, because the other terms can depend on the sampled attention.\par
\#\#\#\# Counterevidence and alternative readings.\par
The words “negative log-likelihood” suggest that the authors may have intended Eq. [eq:nll] (\#eq:nll). Enlarging the logarithm’s parentheses alone would still give a log-likelihood with the wrong sign for minimization. A full correction must change both sign and scope. A variance output \(\sigma^2\) versus standard-deviation output \(\sigma\) does not resolve this issue. A bounded positive scale does not restore the NLL identity either, though it can remove the unbounded limit. No bound is stated in the inspected specification. Reported positive ablation results are compatible with a correct implementation despite an incorrect equation. With no implementation in the packet, that possibility remains unresolved.\par
\#\#\#\# Impact and repair scope.\par
The assessment is Major for the invalid likelihood-based specification of the core uncertainty module: Eq. (6) defines its regularizer and Eq. (10) directly consumes it. This is not an assertion that the reported tracker fails. Replacing the summand by Eq. [eq:nll] (\#eq:nll) restores the Gaussian negative-log-density identity, by taking the logarithm explicitly. It does not establish calibrated predictive uncertainty or reproduce the 4.4-point UPD gain. Moreover, the corrected NLL with freely learned scale can itself collapse at zero residual; practical scale constraints and the intended statistical target must be specified. No complete algorithmic repair or author adoption has been verified.\par
\# Other checks and surviving evidence\par
Table 2’s headline difference is arithmetically correct: \(100(0.663-0.574)=8.9\) percentage points. It compares ViT-L with V2-99, explicitly shown in the table. The same-backbone V2-99 difference is 3.4 points; the larger headline difference cannot isolate the three modules from backbone effects. The controlled small-resolution ablations in Table 3 give UPD, PQI and HQD improvements of 4.4, 2.4 and 2.9 points, respectively, and a combined 6.4 points. Those are internally consistent and need not be additive.\par
The denoising partition is well-defined for \(0\leq\beta_{\rm lower}<\beta_{\rm upper}\leq1\): positive and negative strict inequalities are disjoint and intermediate values are ignored. Equality at a threshold can naturally join the ignored set, but the implementation is not available to check it. The backprojection formula in Section 3.2 is consistent with pinhole depth backprojection when depth means optical-axis depth and the rigid transform includes the usual homogeneous coordinate. These conditional observations do not validate the detector or its training.\par
Figures 4, 6 and 7 were viewed as source images. They illustrate selected occlusion examples, initial query locations, boxes and attention concentration; they do not supply calibration curves or a measured coverage guarantee for predictive uncertainty. Appendix B.2 offers a qualitative rationale, while Table 3 reports entropy and standard deviation. Lower uncertainty alone is not proof of better calibration. This is a limit of the evidence, not a second demonstrated error. Raw predictions, training seeds, evaluator configuration and the unavailable video were not inspected. No issue is asserted about the benchmark arithmetic or the PQI/HQD modules beyond the stated scope.\par
\# Recorded state\par
Dimension | State\par
Inspection assessment | Major\par
Scientific adjudication | Pending\par
Repair verification | Limited to statements explicitly proved here\par
Author adoption | Not established\par
Human reviews | None supplied\par
9 Tao Tang; Lijun Zhou; Pengkun Hao; Zihang He; Kalok Ho; Shuo Gu; Zhihui Hao; Haiyang Sun; Kun Zhan; Peng Jia; Xianpeng Lang; Xiaodan Liang. S2-Track: A Simple yet Strong Approach for End-to-End 3D Multi-Object Tracking. PMLR 267, 2025. Frozen published version. https://proceedings.mlr.press/v267/tang25p.html.\par
\# Evidence and reproducibility\par
The source SHA-256 is `a737bc6e6fc24ef849562db83ba8109cf8a69708494c22b2f7f25d96e14d629a`. Independent calculations, their exact outputs, input hashes, static code anchors, coverage and build records are stored under `verification/0314/`. The script `independent\_check.py` is reviewer-written and imports no author program. Numerical checks support the explicit bounded arguments above and are not benchmark reproductions. Source illustrations inspected as images are identified in the coverage ledger. All report pages are subject to a separate hash-bound visual record after compilation. AI performed this initial review; no human review or independent acceptance is claimed. Separate OpenReview attachments inaccessible during source preparation remain unavailable. Later revisions and errata beyond the frozen source packet were not independently searched; no novelty or priority of discovery is asserted.\par
\end{document}
